\section{Introduction}
\label{sec:intro}

Deep learning has entered an era dominated by large pre-trained foundation models. Fine-tuning a pre-trained base model on task-specific datasets remains the standard paradigm for achieving state-of-the-art performance across diverse domains. However, serving a dedicated, full-sized neural network expert for every single target task is prohibitively expensive in terms of storage, memory, and serving infrastructure, especially as base models scale to billions of parameters. 

To address this deployment bottleneck, post-hoc weight-space fusion or \emph{model merging} has emerged as an elegant, training-free solution~\cite{wortsman2022model,ilharco2022editing}. Standard model merging, such as \emph{Task Arithmetic}~\cite{ilharco2022editing}, treats the parameters of fine-tuned models as vector offsets (task vectors) relative to the pre-trained weights. By linearly combining these task vectors, one can theoretically fuse multiple distinct capabilities into a single model without any additional training cost or active parameter expansion during inference.

Despite its conceptual appeal, standard model merging suffers from a fundamental limitation: \emph{catastrophic representational collision} (or interference) in weight space~\cite{yadav2023ties}. When independent models are fine-tuned on different tasks, they adapt their parameter configurations in orthogonal directions to minimize their respective losses. Fusing these updates via naive addition or averaging causes severe destructive interference, disrupting the fragile structural representations of the base model and leading to a significant collapse in performance (often reducing accuracies to near-random guessing, as observed in our experiments).

To mitigate this weight-space interference, recent literature has proposed complex heuristics. For instance, TIES-Merging~\cite{yadav2023ties} removes low-magnitude parameter updates, performs sign election to resolve sign conflicts, and averages only sign-compatible parameters. Similarly, DARE-Merging~\cite{yu2023language} randomly drops a high percentage of weights and rescales the remainder to maintain activation expectation. However, these methods introduce considerable complexity, such as stochastic behavior, or rely on multi-stage sign-compatibility checks that may not be necessary. 

In this work, we adopt a strictly \emph{empirical} philosophy: we argue that the primary mechanism for mitigating representational collapse is simple weight-space regularization. Rather than designing complex, multi-stage, non-linear coordinate-warping or unconstrained adaptation pipelines, we propose \textbf{Sparsity-Guided Task Arithmetic (SG-TA)}. SG-TA is a remarkably simple and deterministic weight-space regularization framework that decouples the sparsification phase. It applies direct magnitude-based binary masking to individual task-specific update vectors prior to merging. This masking filters out low-magnitude updates, which act as orthogonal parameter noise, preserving only the essential high-magnitude updates that define task specialization.

Furthermore, we explore two distinct masking scopes to understand how parameter budget is best allocated across the model:
\begin{itemize}
    \item \textbf{Global Quantile (GQ) Masking:} Computes a single magnitude threshold globally across the entire model for each task vector, allowing different layers to retain varying percentages of parameters.
    \item \textbf{Layer-wise Quantile (LQ) Masking:} Computes independent magnitude thresholds for each model layer, enforcing a homogeneous parameter keep-ratio across all layers.
\end{itemize}

To optimize the keep-ratio $k$ and merging scaling coefficient $\alpha$ without overfitting, we leverage \emph{Offline Few-Shot Validation Tuning (OFS-Tune)}~\cite{regcalmerge}, utilizing a tiny, robust validation set of only 10 samples per task. This bypasses the \emph{Overfitting-Optimizer Paradox}~\cite{wang2020tent} commonly observed in unsupervised test-time adaptation.

Adhering to our core empirical values, we do not rely on theoretical abstractions. Instead, we subject SG-TA to exhaustive, multi-axial grid sweeps over keep-ratios $k \in [0.1, 1.0]$, masking granularities, and coefficients $\alpha$ across 4 diverse visual classification domains (\textbf{MNIST}, \textbf{FashionMNIST}, \textbf{CIFAR-10}, and \textbf{SVHN}) using a Vision Transformer (\texttt{vit\_tiny\_patch16\_224}) backbone.

We explicitly acknowledge that while our evaluation is centered on a compact Vision Transformer (\texttt{vit\_tiny\_patch16\_224}, 5.7M parameters) and low-resolution benchmarks, this design serves as a highly controlled and computationally efficient sandbox to isolate and dissect the precise mechanics of weight-space masking. This allows us to perform massive parallel sweeps over seeds, keep-ratios, and baseline parameters that would otherwise be computationally prohibitive on LLM-scale models, establishing clear, statistically rigorous trends that lay the groundwork for future larger-scale verification.

Our empirical findings yield several key insights:
\begin{enumerate}
    \item \textbf{Spatial Regularization Hypothesis Validated:} Restricting the active parameter volume via simple magnitude-based masking acts as an excellent regularizer, drastically reducing weight-space interference and preventing representational collapse.
    \item \textbf{Global Budget Flexibility is Crucial:} Global Quantile (GQ) masking consistently and substantially outperforms Layer-wise Quantile (LQ) masking (achieving an average joint mean accuracy of \textbf{61.40\% $\pm$ 1.39\%} vs. \textbf{57.81\% $\pm$ 2.52\%} at their best configurations). This demonstrates that enforcing homogeneous layer-wise budgets is highly sub-optimal, and allowing key layers (such as attention projections) to retain more updates globally is essential.
    \item \textbf{Superiority Over Complex Baselines:} Simple deterministic GQ masking outperforms complex stochastic methods like DARE-Merging (58.44\% $\pm$ 3.02\%) and multi-stage consensus protocols like TIES-Merging (60.64\% $\pm$ 1.30\%) on our benchmark, showcasing that resolving sign conflicts is of secondary importance compared to optimal magnitude-based regularization. However, we note with scientific honesty that because of overlapping standard deviations, our method's superiority over TIES-Merging is not statistically significant.
    \item \textbf{The Absolute Degradation Bottleneck:} Despite significant relative improvements over unregularized baselines, we observe a substantial absolute performance gap of $34.51\%$ between the merged model ($61.40\%$) and the joint expert ceiling ($95.91\%$). This highlights a severe capacity constraint in compact architectures, rendering the resulting merged model practically undeployable for high-stakes applications.
\end{enumerate}

To the best of our knowledge, this is the most exhaustive empirical study characterizing decoupled, deterministic weight-space masking for model merging on Vision Transformers. We provide reproducible code and detailed sweeps to establish a strong, verified baseline for the community.
